\hypertarget{appendix-ffdc1d4e5350}{}\label{referencestart:appendix-ffdc1d4e5350}
\section{Matrix storage indexing and strides}\label{appendix-ffdc1d4e5350-matrix-storage-indexing-and-strides}

\textbf{Curriculum depth:} Introduced · \textbf{First seen:} Phase7 · \textbf{Latest development:} Phase7

\begin{allmightroledefinition}

\subsection{Definition and mechanism}\label{appendix-ffdc1d4e5350-definition-and-mechanism}

Matrix row/column semantics are independent of physical row-major or column-major storage. A strided view uses explicit row and column strides, shape, backing extent, and borrow lifetime. Separate named scalar members are not an array eligible for pointer indexing.

\end{allmightroledefinition}

\textbf{Reconstruction cue:} For row-major strides row=8 and column=1, element (2,3) has scalar offset 19.

\subsection{Prerequisites}\label{appendix-ffdc1d4e5350-prerequisites}

\begin{itemize}
\tightlist
\item
  \href{../01-Bits-Programs-and-Reliable-Systems/main.pdf\#nameddest=appendix-c4ac6bb0fff8}{ABI and object layout (Bits, Programs, and Reliable Systems, Reference notes: ABI and object layout)}
\item
  \href{../01-Bits-Programs-and-Reliable-Systems/main.pdf\#nameddest=appendix-6af5921bc853}{Pointer arithmetic and provenance (Bits, Programs, and Reliable Systems, Reference notes: Pointer arithmetic and provenance)}
\item
  \hyperref[appendix-bd39c3bf5d64]{Matrices as linear maps and basis images}
\end{itemize}

\subsection{Used by}\label{appendix-ffdc1d4e5350-used-by}

\begin{itemize}
\item
  \hyperref[appendix-1bfb7c657105]{Matrix SIMD tiling and cache reuse}
\item
  \href{../07-Vision-XR-and-Playable-Worlds/main.pdf\#nameddest=appendix-55e4b5b0a1e6}{LU factorization and reusable solves (Vision, XR, and Playable Worlds, Reference notes: LU factorization and reusable solves)} --- Extends this foundation into linear-systems and decomposition analysis.
\item
  \href{../07-Vision-XR-and-Playable-Worlds/main.pdf\#nameddest=appendix-67c87a3d9e94}{Sparse and iterative solvers (Vision, XR, and Playable Worlds, Reference notes: Sparse and iterative solvers)} --- Extends this foundation into linear-systems and decomposition analysis.
\item
  \href{../04-Light-Transport-and-Image-Synthesis/main.pdf\#nameddest=appendix-749ba51a6c25}{GPU vertex specification (Light Transport and Image Synthesis, Reference notes: GPU vertex specification)} \#\# Recall and next depth
\end{itemize}

Reconstruct the example, state its conventions and find an input that violates its preconditions. Deeper algorithms and measured production performance remain separate from the present model.

\subsection{Source and navigation}\label{appendix-ffdc1d4e5350-source-and-navigation}

\hyperref[01-foundations-i-ch018-keep-the-promises-visible-while-moving-between-representations]{Continuous explanation} · \hyperref[appendix-fd090d982a15-phase-7---parts-601-700]{Coverage and source gap} · \hyperref[appendix-a103e7846d21-phase-7---source-analysis-and-verification]{Sources and verification}
